\begin{rSection}{Experience}

\begin{rSubsection}{Keysight Technologies}{May 2026 -- Aug 2026}{Machine Learning Security Intern --- Applied ML \& Fault Injection}{Santa Rosa, CA}
\item Developing model-guided optimization and adaptive-learning workflows that automatically tune fault-injection parameters, targeting higher vulnerability-detection coverage than baseline sweep-based flows.
\item Designing experiments and telemetry instrumentation; building scalable hyperparameter search and reliability metrics; delivering a prototype demonstration on production device-security test tooling.
\end{rSubsection}

\begin{rSubsection}{The University of Kansas}{Jan 2022 -- Present}{Graduate Research Assistant, Hardware Security (Advisor: Dr.\ Tamzidul Hoque)}{Lawrence, KS}
\item Designed Trojan-assisted meta-attacks that systematically locate and neutralize design-for-security primitives in microelectronic designs; demonstrated authentication bypass of ring-oscillator PUFs and unlocking of dynamically obfuscated scan chains (ACM TODAES 2026).
\item Led an in-depth security evaluation of lightweight homomorphic obfuscation schemes that protect software secrets against hardware Trojans (IACR TCHES 2026).
\item Built SPHINX, an automated framework for identifying and extracting security primitives from gate-level netlists --- reverse-engineering tooling that maps a design's defensive surface for attack planning (GLSVLSI 2026).
\item Developed golden-chip-free, ML-assisted power side-channel techniques for detecting and localizing hardware Trojans on FPGAs (Journal of Hardware and Systems Security 2026; IEEE NAECON 2023).
\item Built HOACS, a software-only mitigation that keeps secrets encoded during execution on untrusted COTS processors via residue number coding, implemented as C-library and LLVM-based program transformations for cryptographic and arithmetic workloads.
\item Proposed a gray-box, pre-deployment security-evaluation methodology for opaque COTS hardware, driven by power and EM side-channel measurements under zero-trust supply-chain assumptions.
\item Demonstrated power side-channel and fault attacks on the ASCON lightweight-cryptography standard and evaluated countermeasures; delivered hands-on power-SCA tutorials and workshops.
\item Mentored four NSF REU undergraduate researchers across two universities; co-developed and co-taught a hands-on computer hardware course using FPGAs and microcontrollers (NSF-funded).
\end{rSubsection}

\end{rSection}
